\documentclass[10pt]{resume}

\usepackage{fontawesome} % For icons
\usepackage{hyperref} % For better link handling
\usepackage[left=0.4in,top=0.3in,right=0.4in,bottom=0.3in]{geometry} % Slightly reduced margins

\newcommand{\tab}[1]{\hspace{.2667\textwidth}\rlap{#1}}
\newcommand{\itab}[1]{\hspace{0em}\rlap{#1}}
\name{Maruan Bakri Ottoni}
\address{
    \faMobile\hspace{0.5em}+55 41 99209-2810 \hspace{1em}
    \faEnvelope\hspace{0.5em}\href{mailto:maruanbakriottoni@gmail.com}{maruanbakriottoni@gmail.com} \hspace{1em}
    \faGithub\hspace{0.5em}\href{https://github.com/mbottoni/}{Github} \hspace{1em}
    \faDesktop\hspace{0.5em}\href{https://mbottoni.github.io/}{mbottoni.github.io}
}

\begin{document}
\small

\begin{rSection}{Education}
\textbf{MSc in Computer Science}, University of São Paulo (USP) \hfill \textit{Ongoing} \\
Research focus on computational neuroscience and random graph modeling of neural connections and social networks\\
\textit{Coursework}: Graph Theory, Statistical Methods, Deep Learning, NLP, Parallel Computing

\vspace{0.2em}
\textbf{BSc in Electrical Engineering}, University of Campinas (Unicamp), Brazil \hfill \textit{2018 -- 2022} \\
\hfill GPA: 8.5/10 \\
\textit{MSc/PhD level coursework}: Information Theory, Stochastic Processes, Bayesian Inference, Machine Learning \\
\textit{Research Projects}: Introduction to Lattice-based Cryptography(CNPQ Scholarship 2021-2022), Information Geometry for Machine Learning (CNPQ Scholarship 2022-2023)
\end{rSection}

\begin{rSection}{Professional Experience}
\textbf{Data Scientist II} \hfill September 2023 - Present\\
CloudWalk Inc. \hfill \textit{São Paulo, Brazil} \\
\itemsep -3pt {}
As part of an exploration team initiative, my job was to propose new state-of-the-art ideas to improve AI usage across the company. These include:
\begin{itemize}
    \item Designed and deployed scalable monitoring tools to track the health of over 3000 ML features in our feature store, improving model reliability and observability.
    \item Built generative diffusion models for synthetic data generation in the context of tabular features for downstream tasks such as credit scoring and fraud detection, improving classification metrics by approximately $1\%$.
    \item Pretrained and finetuned foundational models for credit assignment using LLM-inspired architectures, enabling the deployment of credit scoring solutions that led to over R\$ 100 million in increased credit approval for clients.
    \item Developed predictive models for credit release performance, significantly accelerating the experimental phase and streamlining the deployment of new credit initiatives and models.
    \item Engineered risk assessment models for the credit portfolio to monitor and mitigate potential financial losses, ensuring the stability and health of the company's credit offerings.
    \item Led end-to-end development of experimental AI projects involving self-supervised learning, graph-based learning (GNNs) for fraud detection, and mechanistic interpretability for Large Language Models (LLMs).
\end{itemize}
\textbf{Quant Analyst Summer Internship} \hfill January 2023 - March 2023\\
Goldman Sachs \hfill \textit{São Paulo, Brazil} \\
\itemsep -3pt {}
As part of a quant team I helped the development of a linear optimization software for asset inventory management across different entities of the bank
 \begin{itemize}
    \item Engineered algorithmic solutions that optimized asset allocation and significantly reduced redundant asset movements, improving operational efficiency.
\end{itemize}

\textbf{Quant Analyst Internship} \hfill October 2022 - November 2022\\
UHedge Trading Solutions \hfill \textit{Campinas, Brazil}
\begin{itemize}
    \itemsep -3pt {}
    \item Conducted mathematical analysis of derivatives strategies to evaluate
    trading opportunities and risk-return profiles
\end{itemize}
\end{rSection}

\begin{rSection}{Relevant Side Projects}

\textbf{Logit Graph} Developed a library for probabilistic graph generation via a logit-based edge process. Did also iterative benchmarks on current statistical graph models, achieving state of art performance
for the method developed. \href{https://pypi.org/project/logit-graph/}{Lib}

\textbf{Horror Movie Recommendation} Built a FastAPI app that recommends horror movies
from OMDb based on user prompts using a unified retrieval
pipeline: Sentence-Transformers embeddings with
blended priors (IMDb rating/votes, facet overlap, recency), optional
cross-encoder reranking. Added UI, PostgreSQL, CI/CD, Payments integrations.
\href{https://github.com/mbottoni/terror-reco}{Code}

\textbf{RAG Solution for Analyzing Auction Documents} Developed a Retrieval-Augmented Generation (RAG) chatbot using Langchain and MongoDB, with a chat interface built in Flask. The bot enables interaction with potentially hundreds or thousands of documents across multiple auctions, allowing users to efficiently navigate and retrieve relevant information. \href{https://github.com/LLAP-Capital/document-analysis}{Code}

\textbf{Graph Similarity via Representation Learning} Contributed with ML methods for graph distance, implementing SOTA architectures and matching strong baselines. \href{https://github.com/dcuoliveira/graph-corr-embedd}{Code}

\end{rSection}


\begin{rSection}{Skills \& Languages}
\begin{tabular}{ @{} >{\bfseries}l @{\hspace{4ex}} l }
Programming & Python, Matlab, C, Java \\
Tools & Google Cloud Platform, Docker, Git, Linux, Kubernetes, Spark, Airflow \\
ML/AI & PyTorch, TensorFlow, Scikit-learn, FSDP Training, Hugging Face\\
Data/Viz & SQL, Matplotlib, Seaborn, Plotly \\
Languages & Portuguese (Native), English (Fluent C2), French (Basic A2) \\
\end{tabular}
\vspace{-0.4em}
\end{rSection}

\end{document}
